\addchap{Úvod}
\label{chap:uvod}

V posledních letech došlo k zásadnímu posunu v organizaci práce a správě síťové infrastruktury. S masivním přechodem na práci na dálku a distribuované týmy se tradiční modely zabezpečení sítě ukazují jako nedostatečné. Organizace stále častěji adoptují architektury typu Zero Trust Network Access (ZTNA), kde je bezpečné VPN připojení klíčovým prvkem nejen pro koncové uživatele, ale i pro serverovou infrastrukturu.

Platforma GoodAccess v současné době nabízí řešení primárně prostřednictvím grafického uživatelského rozhraní (GUI). Toto řešení je vhodné pro běžné uživatele na osobních počítačích, avšak naráží na limity v prostředí serverů, automatizovaných systémů a DevOps pracovních postupů.

V serverovém prostředí Linux, které je dominantní platformou pro cloudovou infrastrukturu, často není grafické rozhraní dostupné (tzv. headless servery). Administrátoři a DevOps inženýři vyžadují nástroje, které lze ovládat z příkazové řádky, snadno skriptovat a integrovat do CI/CD pipelin. Absence nativního CLI (Command Line Interface) klienta pro Linux představuje pro GoodAccess značné omezení v možnosti nasazení v těchto scénářích.
